\chapter{Chapter05 — Gaussian Sharp Level One Counterexample}
%%% generated by the blueprint pipeline from FABL.Chapter05.GaussianSharpLevelOneCounterexample
\section{Overview}

Declarations of the Lean module \texttt{FABL.Chapter05.GaussianSharpLevelOneCounterexample}.

\section{Declarations}

\begin{definition}[Signed Hamming two-tail function]
\lean{FABL.signedHammingTwoTail}
\label{FABL.signedHammingTwoTail}
\leanok
\uses{FABL.hammingUpperTailIndicator}
For a natural number $n$, a real threshold $t$, and a sign vector $x\in\{-1,1\}^n$, let
$H_{t,n}$ denote the Hamming upper-tail indicator on the sign cube. The signed two-tail
function is $H_{t,n}(x)-H_{t,n}(-x)$, where $-x$ is the coordinatewise negation of $x$.
\end{definition}

\begin{theorem}[Oddness of the normalized sign sum]
\lean{FABL.normalizedRademacherSum_neg_counterexample}
\label{FABL.normalizedRademacherSum_neg_counterexample}
\leanok
\uses{FABL.Sign, FABL.linearForm, FABL.normalizedRademacherSum, FABL.signValue}
\difficulty{4}
For every natural number $n$ and every sign vector $x\in\{−1,1\}^n$, the normalized sum
of its coordinates changes sign when all coordinates are negated.
\end{theorem}

\begin{theorem}[Monomial under global sign reversal]
\lean{FABL.monomial_neg_input_counterexample}
\label{FABL.monomial_neg_input_counterexample}
\leanok
\uses{FABL.monomial, FABL.signValue}
\difficulty{1}
Let $n$ be a natural number, let $S$ be a subset of the $n$ coordinate indices, and let
$x$ be a vector of $n$ signs. Reversing every sign in $x$ multiplies the monomial
indexed by $S$ by $(-1)^{|S|}$, where $|S|$ is the number of indices in $S$.
\end{theorem}

\begin{theorem}[Sign reversal of a singleton Fourier coefficient]
\lean{FABL.fourierCoeff_comp_neg_singleton}
\label{FABL.fourierCoeff_comp_neg_singleton}
\leanok
\uses{FABL.fourierCoeff, FABL.monomial, FABL.monomial_neg_input_counterexample}
\difficulty{4}
For every natural number $n$, every real-valued function $f$ on the $n$-dimensional sign
cube, and every coordinate $i$, the Fourier coefficient at $\{i\}$ of the function
$x\mapsto f(-x)$ is the negative of the Fourier coefficient of $f$ at $\{i\}$.
\end{theorem}

\begin{theorem}[Singleton coefficient of a signed tail]
\lean{FABL.fourierCoeff_signedHammingTwoTail_singleton}
\label{FABL.fourierCoeff_signedHammingTwoTail_singleton}
\leanok
\uses{FABL.fourierCoeff, FABL.fourierCoeff_comp_neg_singleton, FABL.hammingUpperTailIndicator, FABL.monomial, FABL.signedHammingTwoTail}
\difficulty{4}
Let $n\geq 1$, let $t\in\mathbb R$, and choose a coordinate $i\in[n]$. The Fourier
coefficient at $\{i\}$ of the signed two-sided tail of the normalized Rademacher sum
equals twice the Fourier coefficient at $\{i\}$ of its upper-tail indicator.
\end{theorem}

\begin{theorem}[Absolute value of the signed two-tail function]
\lean{FABL.abs_signedHammingTwoTail}
\label{FABL.abs_signedHammingTwoTail}
\leanok
\uses{FABL.hammingUpperTailIndicator, FABL.normalizedRademacherSum, FABL.normalizedRademacherSum_neg_counterexample, FABL.signedHammingTwoTail}
\difficulty{4}
Let $n$ be a natural number, let $x\in\{\pm1\}^n$, and let $t\ge 0$. The absolute value
of the difference between the upper-tail indicators at $x$ and at $-x$ equals their sum.
\end{theorem}

\begin{theorem}[Signed two-sided tail bounds]
\lean{FABL.signedHammingTwoTail_mem_Icc}
\label{FABL.signedHammingTwoTail_mem_Icc}
\leanok
\uses{FABL.hammingUpperTailIndicator, FABL.signedHammingTwoTail}
\difficulty{4}
For every real threshold $t$, natural number $n$, and sign vector $x\in\{-1,1\}^n$, the
difference between the upper-tail indicator at $x$ and the upper-tail indicator at $-x$
lies in $[-1,1]$. Here the upper-tail indicator is one when the normalized Rademacher
sum is strictly greater than $t$, and zero otherwise.
\end{theorem}

\begin{theorem}[Symmetry of a Rademacher tail]
\lean{FABL.expect_hammingUpperTailIndicator_comp_neg}
\label{FABL.expect_hammingUpperTailIndicator_comp_neg}
\leanok
\uses{FABL.hammingUpperTailIndicator}
\difficulty{3}
For every $t\in\mathbb R$ and $n\in\mathbb N$, the uniform expectation over the sign
cube of the indicator that the normalized Rademacher sum of $-x$ exceeds $t$ equals the
uniform expectation of the same indicator at $x$.
\end{theorem}

\begin{theorem}[Absolute signed two-tail expectation]
\lean{FABL.expect_abs_signedHammingTwoTail}
\label{FABL.expect_abs_signedHammingTwoTail}
\leanok
\uses{FABL.abs_signedHammingTwoTail, FABL.expect_hammingUpperTailIndicator_comp_neg, FABL.hammingUpperTailIndicator, FABL.signedHammingTwoTail}
\difficulty{4}
Let $t\geq0$ and let $n$ be a nonnegative integer. For a uniformly random sign vector
$X\in\{-1,1\}^n$, define $R(X)=n^{-1/2}\sum_{i=1}^n X_i$ and
$U_t(X)=\mathbf{1}_{\{t<R(X)\}}$. Then $\mathbb{E}|U_t(X)-U_t(-X)|=2\mathbb{E}U_t(X)$.
\end{theorem}

\begin{theorem}[Signed Tail Fourier Weight]
\lean{FABL.fourierWeightAtLevel_one_signedHammingTwoTail}
\label{FABL.fourierWeightAtLevel_one_signedHammingTwoTail}
\leanok
\uses{FABL.fourierCoeff_signedHammingTwoTail_singleton, FABL.fourierWeightAtLevel, FABL.fourierWeightAtLevel_one_eq_sum_singleton, FABL.hammingUpperTailIndicator, FABL.signedHammingTwoTail}
\difficulty{4}
For every real threshold $t$ and natural number $n$, define the signed tail function on
the $n$-dimensional sign cube as the indicator that the normalized coordinate sum
exceeds $t$, minus the same indicator evaluated at the negated input. Its level-one
Fourier weight equals four times the level-one Fourier weight of the upper-tail
indicator.
\end{theorem}

\begin{definition}[Signed Hamming tail coefficient]
\lean{FABL.signedHammingTwoTailEpsilon}
\label{FABL.signedHammingTwoTailEpsilon}
\leanok
\uses{FABL.fourierCoeff, FABL.hammingUpperTailIndicator}
For a real threshold $t$ and a natural number $m$, define the signed Hamming tail
parameter to be the absolute value of twice the Fourier coefficient of the Hamming
upper-tail indicator on the $(m+1)$-dimensional sign cube at the singleton set
containing coordinate $0$.
\end{definition}

\begin{theorem}[Singleton coefficient of a signed tail function]
\lean{FABL.abs_fourierCoeff_signedHammingTwoTail_eq_epsilon}
\label{FABL.abs_fourierCoeff_signedHammingTwoTail_eq_epsilon}
\leanok
\uses{FABL.fourierCoeff, FABL.fourierCoeff_hammingUpperTailIndicator_singleton, FABL.fourierCoeff_signedHammingTwoTail_singleton, FABL.signedHammingTwoTail, FABL.signedHammingTwoTailEpsilon}
\difficulty{3}
Let $t$ be real, let $m$ be a nonnegative integer, and let $i$ be a coordinate in
dimension $m+1$. The absolute value of the Fourier coefficient in coordinate $i$ of the
signed two-sided Hamming tail function at threshold $t$ equals the signed-tail parameter
defined by twice the absolute first-level Fourier coefficient of the upper-tail
indicator in dimension $m+1$.
\end{theorem}

\begin{theorem}[Reciprocal square root limit]
\lean{FABL.tendsto_inv_sqrt_nat_succ}
\label{FABL.tendsto_inv_sqrt_nat_succ}
\leanok
\difficulty{3}
As the natural number $m$ tends to infinity, the reciprocal of the square root of $m+1$
tends to $0$.
\end{theorem}

\begin{theorem}[Vanishing signed Hamming tail coefficient]
\lean{FABL.tendsto_signedHammingTwoTailEpsilon}
\label{FABL.tendsto_signedHammingTwoTailEpsilon}
\leanok
\uses{FABL.fourierCoeff, FABL.fourierCoeff_hammingUpperTailIndicator_singleton, FABL.hammingUpperTailIndicator, FABL.signedHammingTwoTailEpsilon, FABL.tendsto_expect_hammingUpperTailIndicator_mul_normalizedRademacherSum, FABL.tendsto_inv_sqrt_nat_succ}
\difficulty{4}
For every real number $t$, consider the uniform sign cube of dimension $m+1$, the sum of
its coordinates divided by $\sqrt{m+1}$, and the indicator of the event that this
normalized sum exceeds $t$. Twice the absolute value of the Fourier coefficient of this
indicator at the singleton consisting of the first coordinate tends to $0$ as
$m\to\infty$.
\end{theorem}

\begin{theorem}[Limit of the signed Hamming tail expectation]
\lean{FABL.tendsto_expect_abs_signedHammingTwoTail}
\label{FABL.tendsto_expect_abs_signedHammingTwoTail}
\leanok
\uses{FABL.expect_abs_signedHammingTwoTail, FABL.hammingUpperTailIndicator, FABL.signedHammingTwoTail, FABL.standardGaussianUpperTail, FABL.tendsto_expect_hammingUpperTailIndicator}
\difficulty{3}
For every real $t\ge0$, the uniform-cube expectation in dimension $m+1$ of the absolute
value of the difference between the upper-tail indicators of the normalized Rademacher
sum at $x$ and at $-x$ converges, as $m\to\infty$, to twice the standard Gaussian
upper-tail probability at $t$.
\end{theorem}

\begin{theorem}[First-level Fourier weight of a signed tail function]
\lean{FABL.tendsto_fourierWeightAtLevel_one_signedHammingTwoTail}
\label{FABL.tendsto_fourierWeightAtLevel_one_signedHammingTwoTail}
\leanok
\uses{FABL.fourierWeightAtLevel, FABL.fourierWeightAtLevel_one_signedHammingTwoTail, FABL.hammingUpperTailIndicator, FABL.proposition5_25, FABL.signedHammingTwoTail}
\difficulty{3}
For each real number $t$, let $F_n$ be the function on the sign cube $\\{-1,1\\}^n$ that
equals the indicator of $n^{-1/2}\\sum_{i=1}^n x_i>t$ minus the indicator of
$n^{-1/2}\\sum_{i=1}^n (-x_i)>t$. As $m$ tends to infinity, the sum of the squared
Fourier coefficients of $F_{m+1}$ over all singleton coordinate sets tends to
$4\\,\\varphi(t)^2$, where $\\varphi$ is the standard Gaussian density.
\end{theorem}

\begin{theorem}[Strict concavity of the Gaussian isoperimetric profile]
\lean{FABL.strictConcaveOn_gaussianIsoperimetricReal}
\label{FABL.strictConcaveOn_gaussianIsoperimetricReal}
\leanok
\uses{FABL.continuousOn_gaussianIsoperimetricReal, FABL.gaussianIsoperimetricReal, FABL.gaussianIsoperimetricReal_pos, FABL.gaussianIsoperimetric_second_derivative}
\difficulty{3}
The real-valued Gaussian isoperimetric profile is strictly concave on $[0,1]$.
\end{theorem}

\begin{theorem}[A Gaussian profile inequality in a tail interval]
\lean{FABL.exists_signedHammingTwoTail_alpha}
\label{FABL.exists_signedHammingTwoTail_alpha}
\leanok
\uses{FABL.continuousOn_gaussianIsoperimetricReal, FABL.gaussianIsoperimetricReal, FABL.strictConcaveOn_gaussianIsoperimetricReal}
\difficulty{5}
There exists $\alpha\in(1/4,1/2)$ such that the standard Gaussian isoperimetric profile
at $\alpha$ is less than twice its value at $1/8$.
\end{theorem}

\begin{theorem}[Exercise5 44 printed false]
\lean{FABL.exercise5_44_printed_false}
\statementsource{Analysis-of-Boolean-Functions-by-Ryan-ODonnell}{pdf-ffbe9c7b3306-p139-b007}
\label{FABL.exercise5_44_printed_false}
\leanok
\uses{FABL.abs_fourierCoeff_signedHammingTwoTail_eq_epsilon, FABL.exists_signedHammingTwoTail_alpha, FABL.fourierCoeff, FABL.fourierWeightAtLevel, FABL.gaussianIsoperimetric, FABL.gaussianIsoperimetricReal, FABL.gaussianIsoperimetricReal_apply_of_mem_Ioo, FABL.gaussianIsoperimetricReal_eq_gaussianIsoperimetric, FABL.gaussianIsoperimetricReal_pos, FABL.signedHammingTwoTail, FABL.signedHammingTwoTailEpsilon, FABL.signedHammingTwoTail_mem_Icc, FABL.standardGaussianUpperQuantile, FABL.standardGaussianUpperTail, FABL.standardGaussianUpperTail_neg, FABL.standardGaussianUpperTail_quantile, FABL.standardGaussianUpperTail_strictAnti, FABL.tendsto_expect_abs_signedHammingTwoTail, FABL.tendsto_fourierWeightAtLevel_one_signedHammingTwoTail, FABL.tendsto_signedHammingTwoTailEpsilon}
\difficulty{6}
% dataset-pipeline: stub -- prose pending build_informal_statements.py
Statement pending regeneration.
\end{theorem}
